\section{Wreath-Product Farahat--Higman Algebras} \label{rgamma_section}

We now extend the results of Section \ref{FH_theory_section} from the symmetric group $S_n$ to the wreath products $\Gamma \wr S_n = \Gamma ^n \rtimes S_n$, where $\Gamma$ is a finite group. There is a wreath-product version of the Farahat--Higman algebra, which is also related to the ring of symmetric functions. The homomorphisms $\Phi_n$, $ev_n$, $\Psi$ all generalise to the wreath-product setting. We use the same notation for the wreath-product versions of these maps because taking $\Gamma$ to be the trivial group (so that $\Gamma \wr S_n = S_n$), we recover the maps from Section~\ref{FH_theory_section}.

Recall that the conjugacy classes of $\Gamma \wr S_n$ correspond to multipartitions of total size $n$ indexed by $\Gamma_*$. We use the boldface Greek letters $\bs \lambda, \bs \mu, \bs \nu$ to indicate multipartitions. Thus $\bs{\mu}(c)$ means the partition in $\bs \mu$ indexed by $c \in \Gamma_*$. Similarly to symmetric groups, we have a notion of reduced cycle type. 
\begin{definition}[{\cite[Subsection 2.3]{Wang1}}]
The \emph{partially-reduced cycle type} of an element of $\Gamma \wr S_n$ of cycle type $\bs{\mu}$ is the multipartition obtained by subtracting $1$ from each part of $\bs \mu(1)$, and ignoring any resulting parts of size zero. 
\end{definition}

This is the same as the case of symmetric groups, but applied only to the partition labelled by the identity conjugacy class. The group $\Gamma \wr S_n$ contains an element of partially-reduced cycle type $\bs \mu$ if and only if $n \geq |\bs \mu| + l(\bs \mu(1))$. Later in this section we will need to introduce the notion of \emph{fully-reduced cycle type} which will involve reducing each partition in $\bs \mu$ (this is why we do not use the terminology of ``modified type'' from \cite{Wang1}).

The centre of the integral group ring of $\Gamma \wr S_n$ has a $\mathbb{Z}$-basis consisting of conjugacy class sums. We let $X_{\bs \mu}$ be the sum of all elements of partially-reduced cycle type $\bs \mu$ (which if zero if there are no such elements). Then we have a direct analogue of Theorem \ref{thm:FH_alg_polys}.

\begin{theorem}[{\cite[Theorem 2.13]{Wang1}}] \label{FH_g_alg_polys_thm}
For any multipartitions $\bs \lambda, \bs\mu, \bs\nu$, there exists a unique integer-valued polynomial $\phi_{\bs\mu, \bs\nu}^{\bs\lambda}(t)$ such that the equation
\[
X_{\bs\mu} X_{\bs\nu} = \sum_{\bs\lambda} \phi_{\bs\mu, \bs\nu}^{\bs\lambda}(n)X_{\bs\lambda}
\]
holds in $Z(\mathbb{Z} \Gamma \wr S_n)$ for all $n \in \mathbb{Z}_{\geq 0}$.
\end{theorem}
This allows us to define an algebra analogous to $\FH$, but for wreath products.

\begin{definition} 
Let $\FHG$ be the free $\mathcal{R}$-module with basis $K_{\bs \mu}$ indexed by all multipartitions $\bs \mu \in \mathcal{P}(\Gamma_*)$. Define a bilinear multiplication on $\FHG$ by
\[
K_{\bs \mu} K_{\bs \nu} = \sum_{\lambda} \phi_{\bs\mu, \bs\nu}^{\bs\lambda}(t) K_{\bs \lambda}
\]
where $\phi_{\bs\mu, \bs\nu}^{\bs\lambda}(t)$ are the polynomials from Theorem \ref{FH_g_alg_polys_thm}, viewed as elements of the base ring $\mathcal{R}$.
\end{definition}

As in the case of $\FH$, $\FHG$ is a commutative, associative, unital $\mathcal{R}$-algebra, and has specialisation homomorphisms.
\begin{theorem} [{\cite[Section 2.5]{Wang1}}] \label{thm:FH_g_spec_hom}
For each $n \in \mathbb{Z}_{\geq 0}$ there is a surjective ring homomorphism $\Phi_n: \FHG \to Z(\mathbb{Z}\Gamma \wr S_n)$ defined by
\[
\Phi_n \left( \sum_{\bs \mu} a_{\bs \mu}(t) K_{\bs\mu} \right) = \sum_{\bs \mu} a_{\bs \mu}(n) X_{\bs \mu},
\]
where $a_{\bs\mu}(t) \in \mathcal{R}$, and $X_{\bs\mu}$ is the sum of all elements of $\Gamma \wr S_n$ of partially-reduced cycle type $\mu$ \textup{(}or zero, if there are no such elements\textup{)}.
\end{theorem}

As we now show, the family of homomorphisms $\Phi_n$ separates elements of $\FHG$. So to prove an identity in $\FHG$, it is enough to pass to $Z(\mathbb{Z}\Gamma \wr S_n)$ by applying $\Phi_n$.
\begin{proposition} \label{specialisation_distinguishing_prop}
Suppose $x,y \in \FHG$ satisfy $\Phi_n(x) = \Phi_n(y)$ for all sufficiently large positive integers $n$. Then $x=y$.
\end{proposition}
\begin{proof}
Consider an element
\[
z = \sum_{\bs \mu} a_{\bs \mu}(t) K_{\bs\mu}.
\]
Then
\[
\Phi_n \left(z \right) = \sum_{\bs \mu} a_{\bs \mu}(n) X_{\bs \mu},
\]
where $X_{\bs \mu}$ is nonzero provided $n \geq |\bs \mu| + l(\bs \mu(1))$, and such nonzero elements form a basis of $Z(\mathbb{Z}\Gamma \wr S_n)$. This means that $\Phi_n(z)$ determines $a_{\bs \mu}(n)$ for all $n$ sufficiently large, which is a Zariski-dense subset of $\mathbb{Z}$. Hence $\Phi_n(z)$ determines each coefficient $a_{\bs \mu}(t) \in \mathcal{R}$, and so $\Phi_n(z)$ determines $z$ itself. 
\end{proof}

We have introduced $\FHG$ as a $\mathcal{R}$-algebra. However, it will turn out that it is better to view it as an algebra over a $\Gamma_*$-weighted version of $\mathcal{R}$, which we now work towards defining.
\begin{definition}
Consider formal power series in variables $x_c$ indexed by $c \in \Gamma_*$. We encode monomials in the $x_c$ using functions $\mathbf{N}: \Gamma_* \to \mathbb{Z}_{\geq 0}$ that record the exponent of each variable:
\[
x^{\mathbf{N}} = \prod_{c \in \Gamma_*} x_c^{\mathbf{N}(c)}.
\]
Additionally, let $|\mathbf{N}| = \sum_c \mathbf{N}(c)$, so that $\deg(x^{\mathbf{N}}) = |\mathbf{N}|$.
\end{definition}

Note that the set of such functions $\mathbf{N}$ is $\mathbb{Z}_{\geq 0}^{\Gamma_*}$. 
\begin{definition}
Let $\mathfrak{g} = \mathbb{Q}\Gamma_* = Z(\mathbb{Q}\Gamma)$, viewed as a Lie algebra via the associative multiplication \textup{(}it is abelian\textup{)}. There is a canonical map $\iota: \mathfrak{g} \to \mathcal{U}(\mathfrak{g})$ from $\mathfrak{g}$ to its universal enveloping algebra, which we extend to formal power series coefficientwise. We define
\[
T: \mathfrak{g}[[x_{c_1}, \ldots, x_{c_l}]] \to \mathcal{U}(\mathfrak{g})[[x_{c_1},  \ldots, x_{c_l}]]
\]
via
\[
T\left(
\sum_{\mathbf{N} \in \mathbb{Z}_{\geq 0}^{\Gamma_*}} u_{\mathbf{N}} x^{\mathbf{N}}
\right) = \sum_{\mathbf{N} \in \mathbb{Z}_{\geq 0}^{\Gamma_*}} \iota(u_{\mathbf{N}}) x^{\mathbf{N}}.
\]
where $u_{\mathbf{N}} \in \mathfrak{g}$.
\end{definition}
\begin{remark}
The map $T$ allows us to distinguish between the associative multiplication in $\mathfrak{g}$, and the multiplication in the universal enveloping algebra. For example, if $u, v \in \mathfrak{g}$, then $uv \in \mathfrak{g}$, while $T(u)T(v)$ is an element of $\mathcal{U}(\mathfrak{g})$ that lies in degree two with respect to the Poincar\'{e}-Birkhoff-Witt (PBW) filtration. It may seem peculiar to consider the universal enveloping algebra of an abelian Lie algebra when the result is simply the symmetric algebra of the Lie algebra. The ring we are about to define, $\RGamma$, is actually a distribution algebra, and so the universal enveloping algebra is actually the natural place to construct it. This is discussed further in the appendix. There is also similar situation when interpolating Grothendieck rings of wreath products (rather than centres of group algebras of wreath products) where $\mathfrak{g}$ is the Grothendieck ring of a tensor category. In that case, $\mathfrak{g}$ is not necessarily abelian and the universal enveloping algebra is the correct construction. The interested reader is directed to \cite{Ryba1} and \cite{Ryba2}.
\end{remark}
\begin{definition} \label{rgamma_constr_defn}
Let $\Omega \in \mathcal{U}(\mathfrak{g})[[x_{c_1}, \ldots, x_{c_l}]]$ be the generating function defined as follows,
\[
\Omega(x_{c_1}, \ldots, x_{c_l}) = \exp\left(T\left(\log\left(1 + \sum_{c \in \Gamma_*} c x_c\right)\right)\right).
\]
Here $\log\left(1 + \sum_{c \in \Gamma_*} c x_c\right)$ is viewed as an element of $\mathfrak{g}[[x_{c_1}, x_{c_2}, \ldots, x_{c_l}]]$ \textup{(}note that expanding the power series involves the associative algebra structure of $\mathfrak{g}$, not just the Lie algebra structure\textup{)}. Writing $\Omega$ in terms of monomials, we let $B_{\mathbf{N}} \in U(\mathfrak{g})$ be the coefficients:
\[
\Omega(x_{c_1}, \ldots, x_{c_l}) = \sum_{\mathbf{N} \in \mathbb{Z}_{\geq 0}^{\Gamma_*}} B_{\mathbf{N}} x^{\mathbf{N}}.
\]
We let $\RGamma$ be the $\mathbb{Z}$-submodule of $\mathcal{U}(\mathfrak{g})$ spanned by all the $B_{\mathbf{N}}$, and we call $\RGamma$ \emph{the ring of $\Gamma_*$-weighted integer-valued polynomials}.
\end{definition}

We will shortly show that $\RGamma$ is indeed a ring, but first, the following example justifies the analogy to integer-valued polynomials.
\begin{example} \label{trivial_rg_example}
Suppose that $\Gamma$ is the trivial group, so that there is only one conjugacy class $c$, and it obeys $c^2 = c$. Then,
\begin{eqnarray*}
\log\left(1 + \sum_{c \in \Gamma_*} c x_c\right) &=& 
\sum_{i\geq 1} \frac{(-1)^{i-1}}{i} c^i x_c^i \\
&=& c \sum_{i\geq 1} \frac{(-1)^{i-1}}{i} x_c^i \\
&=& c \log(1+x_c).
\end{eqnarray*}
Hence $\Omega$ becomes
\begin{eqnarray*}
\Omega(x_{c}) &=& \exp(T(c \log(1+x_c))) \\
&=& \exp(\log(1+x_c) T(c)) \\
&=& (1+x_c)^{T(c)} \\
&=& \sum_{N \geq 0} \binom{T(c)}{N} x_c^N.
\end{eqnarray*}
So the $B_{\mathbf{N}}$ are simply binomial coefficients (with parameter $T(c)$) and $\RGamma = \mathcal{R}$.
\end{example}

\begin{proposition} \label{theta_product_prop}
We have the following relation:
\[
\Omega(x_{c_1}, \ldots, x_{c_l}) \Omega(y_{c_1}, \ldots, y_{c_l}) = \Omega(z_{c_1}, \ldots, z_{c_l}),
\]
where
\[
z_{c_i} = x_{c_i} + y_{c_i} + \sum_{j,k}A_{j,k}^{i}x_{c_j} y_{c_k}.
\]
\end{proposition}
\begin{proof}
Let $S(x) = 1 + \sum_{c \in \Gamma_*} c x_c$, and similarly define $S(y)$ and $S(z)$. Then
\begin{eqnarray*}
S(x)S(y) &=& \left(1 + \sum_{c \in \Gamma_*} c x_c\right)\left(1 + \sum_{c \in \Gamma_*} c y_c\right) \\
&=& 1 + \sum_{c_j \in \Gamma_*} c_j x_{c_j} + \sum_{c_k \in \Gamma_*} c_k y_{c_k} + \sum_{c_i} A_{j,k}^i c_i x_{c_j}y_{c_k} \\
&=& S(z).
\end{eqnarray*}
Because we are working with commutative algebras, we may compute as follows:
\begin{eqnarray*}
\Omega(x_{c_1}, \ldots, x_{c_l}) \Omega(y_{c_1}, \ldots, y_{c_l})
&=& \exp\left(T\left(\log\left(S(x)\right)\right)\right) \exp\left(T\left(\log\left(S(y)\right)\right)\right) \\
&=& \exp\left(T\left(\log\left(S(x)\right)\right) + T\left(\log\left(S(y)\right)\right)\right) \\
&=&\exp\left(T\left(\log\left(S(x)\right)+\log\left(S(y)\right)\right)\right) \\
&=&\exp\left(T\left(\log\left(S(x)S(y)\right)\right)\right) \\
&=&\exp\left(T\left(\log\left(S(z)\right)\right)\right) \\
&=& \Omega(z_{c_1}, \ldots, z_{c_l}).\qedhere
\end{eqnarray*}
\end{proof}

\begin{lemma} \label{leading_term_lemma}
As an element of $\mathcal{U}(\mathfrak{g})$, $B_{\mathbf{N}}$ lies in PBW filtration degree $|\mathbf{N}|$ and has leading order term
\[
\prod_{c \in \Gamma_*} \frac{T(c)^{\mathbf{N}(c)}}{\mathbf{N}(c)!}.
\]
\end{lemma}
\begin{proof}
The $B_{\mathbf{N}}$ are defined by the equation
\[
\Omega(x_{c_1}, \ldots, x_{c_l}) = \exp\left(T\left(\log\left(1 + \sum_{c \in \Gamma_*} c x_c\right)\right)\right) = \sum_{\mathbf{N} \in \mathbb{Z}_{\geq 0}^{\Gamma_*}} B_{\mathbf{N}} x^{\mathbf{N}}.
\]
The expression
\[
T\left(\log\left(1 + \sum_{c \in \Gamma_*} c x_c\right)\right)
\]
is in PBW degree 1. When the logarithm is expanded as a power series in the $x_c$ variables, the lowest order term is
\[
\sum_{c \in \Gamma_*} T(c) x_c.
\]
In particular, there is no constant term. This means that $\Omega(x_{c_1}, \ldots, x_{c_l})$ is a sum of products of terms whose PBW degree is less than or equal to their degree in the $x_c$ variables. Since $B_{\mathbf{N}}$ is the coefficient of $x^{\mathbf{N}}$, it is contained in PBW filtration degree $|\mathbf{N}|$. Moreover, to compute the leading term of the $B_N$, we neglect all but the lowest degree monomials in $x_c$; to leading order, $\Omega$ is approximated by
\[
\exp\left( \sum_{c \in \Gamma_*} T(c) x_c\right) = 
\prod_{c \in \Gamma_*} \exp\left( T(c) x_c\right).
\]
The leading term of $B_{\mathbf{N}}$ is found by taking the coefficient of $x^{\mathbf{N}}$.
\end{proof}


\begin{proposition} \label{r_gamma_is_an_algebra_prop}
The $B_{\mathbf{N}}$ are a $\mathbb{Z}$-basis of $\RGamma$. Additionally, the multiplication in $\mathcal{U}(\mathfrak{g})$ induces a multiplication on $\RGamma$, making it into a unital commutative ring.
\end{proposition}
\begin{proof}
The $B_{\mathbf{N}}$ span $\RGamma$ by definition. It follows from Lemma \ref{leading_term_lemma} and the PBW theorem that the $B_{\mathbf{N}}$ form a $\mathbb{Q}$-basis of $\mathcal{U}(\mathfrak{g})$, so they are also linearly independent.

We note that when $\mathbf{N}$ is zero, $B_{\mathbf{N}}$ is the constant term of $\Omega$, namely the identity in $\mathcal{U}(\mathfrak{g})$, so $\RGamma$ has an identity element. To show that $\RGamma$ is a commutative ring, we must show that the product of two basis elements is a linear combination of basis elements. This amounts to showing that the coefficient of any monomial $x^{\mathbf{N}} y^{\mathbf{M}}$ in 
\[
\Omega(x_{c_1}, \ldots, x_{c_l}) \Omega(y_{c_1}, \ldots, y_{c_l})
\]
is contained in $\RGamma$. We invoke Proposition \ref{theta_product_prop}, which expresses this product as 
\[
\Omega(z_{c_1}, \ldots, z_{c_l}) = \sum_{\mathbf{K} \in \mathbb{Z}_{\geq 0}^{\Gamma_*}} B_{\mathbf{K}} z^{\mathbf{K}},
\]
where $z_{c_i} = x_{c_i} + y_{c_i} + \sum_{j,k}A_{j, k}^{i}x_{c_j} y_{c_k}$, which is a $\mathbb{Z}$-linear combination of monomials in the $x$ and $y$ variables. Thus the expansion in terms of $x^{\mathbf{N}} y^{\mathbf{M}}$ will have $\mathbb{Z}$-linear combinations of the $B_{\mathbf{K}}$ as coefficients.
\end{proof}

\begin{proposition} \label{rgamma_free_over_r_prop}
The set of $B_{\mathbf{N}}$ in $\RGamma$ indexed by $\mathbf{N}$ with $\mathbf{N}(c) = 0$ for $c \neq 1$ forms a subring of $\RGamma$ isomorphic to $\mathcal{R}$. As a module over $\mathcal{R}$, $\RGamma$ is free with basis $B_{\mathbf{M}}$ indexed by $\mathbf{M}$ with $\mathbf{M}(1) = 0$.
\end{proposition}

\begin{proof}
The $B_{\mathbf{N}}$ with $\mathbf{N}(c) = 0$ for $c \neq 1$ can be found by considering the generating function $\Omega$ and setting all $x_c$ to zero other than $x_1$. Then the calculation in Example \ref{trivial_rg_example} shows $B_{\mathbf{N}} = \binom{T(1)}{\mathbf{N}(1)}$, and binomial coefficients are a $\mathbb{Z}$-basis of $\mathcal{R}$.

To show that $\RGamma$ is free over $\mathcal{R}$ with the stated basis, we pass to the associated graded algebra with respect to the PBW filtration. Lemma \ref{leading_term_lemma} showed the leading order term of $B_{\mathbf{N}}$ is 
\[
\prod_{c \in \Gamma_*} \frac{T(c)^{\mathbf{N}(c)}}{\mathbf{N}(c)!} = 
\frac{T(1)^{\mathbf{N}(1)}}{\mathbf{N}(1)!}  \prod_{c \neq 1} \frac{T(c)^{\mathbf{N}(c)}}{\mathbf{N}(c)!},
\]
which we recognise as the leading term of $\binom{T(1)}{\mathbf{N}(1)} \in \mathcal{R}$ multiplied by the leading term of $B_{\mathbf{M}}$, where $\mathbf{M}(c) = \mathbf{N}(c)$ for $c \neq 1$ and $\mathbf{M}(1) = 0$. This shows that $\RGamma$ is spanned by the $B_{\mathbf{M}}$ over $\mathcal{R}$, and also shows the $\mathcal{R}$-linear independence of the $B_{\mathbf{M}}$.
\end{proof}



\begin{definition}
For any $\mathbf{N} \in \mathbb{Z}_{\geq 0}^{\Gamma_*}$ with $|\mathbf{N}| \leq n$, we define $b_{\mathbf{N}} \in \mathbb{Q}\Gamma \wr S_n$ as follows
\[
b_{\mathbf{N}} = \frac{1}{(n - |\mathbf{N}|)! \prod_{c \in \Gamma_*} \mathbf{N}(c)!}
\sum_{\sigma \in S_n} \sigma \left( 1^{\otimes (n - |\mathbf{N}|)} \otimes \bigotimes_{c \in \Gamma_*} c^{\otimes \mathbf{N}(c)} \right) \sigma^{-1}.
\]
The parenthesised tensor product is an element of $(\mathbb{Q} \Gamma)^{\otimes n} = \mathbb{Q}\Gamma^n \subseteq \mathbb{Q}\Gamma \wr S_n$ \textup{(}because of the averaging over $S_n$, the choice of order of the tensor factors does not matter\textup{)}.
\end{definition}
\begin{lemma}
The element $b_{\mathbf{N}}$ is actually contained in $Z(\mathbb{Z}\Gamma \wr S_n)$.
\end{lemma}
\begin{proof}
Since $b_{\mathbf{N}}$ is defined by averaging over $S_n$, it commutes with $S_n \subseteq \Gamma \wr S_n$. As the element being averaged is a tensor product of central elements of $\mathbb{Z}\Gamma$, $b_{\mathbf{N}}$ will commute with $\Gamma^n \subseteq \Gamma \wr S_n$. This shows that $b_{\mathbf{N}}$ is central. To see that $b_{\mathbf{N}}$ has integer coefficients, note that the conjugation action of the subgroup 
\[
S_{n - |\mathbf{N}|} \times \prod_{c \in \Gamma_*} S_{\mathbf{N}(c)} \subseteq S_n
\]
is trivial. Any element of a fixed coset of this subgroup acts the same way, and the size of this coset cancels out the denominator.
\end{proof}

\begin{lemma} \label{specialised_gen_fun_lemma}
We have the following expression for the generating function for the elements $b_{\mathbf{N}}$:
\[
\sum_{|\mathbf{N}| \leq n} b_{\mathbf{N}} x^{\mathbf{N}} = (1 + \sum_c c x_c)^{\otimes n}.
\]
\end{lemma}

\begin{proof}
We observe that $b_{\mathbf{N}}$ is the sum of pure tensors with $\mathbf{N}(c)$ tensor factors set equal to $c$ and all remaining tensor factors equal to $1$. Here we distinguish tensor factors assigned $c=1$ from unassigned tensor factors which also take the value 1. But the coefficient of $x^{\mathbf{N}}$ on the right hand side will be the sum of all pure tensors with $\mathbf{N}(c)$ factors equal to $c$. 
\end{proof}


\begin{theorem} \label{r_gamma_hom_thm}
There is a homomorphism $ev_n^{\mathcal{R}}: \RGamma \to Z(\mathbb{Z} \Gamma \wr S_n)$ defined by $ev_n^{\mathcal{R}}(B_{\mathbf{N}}) = b_{\mathbf{N}}$ if $|\mathbf{N}| \leq n$, and $ev_n^{\mathcal{R}}(b_{\mathbf{N}}) = 0$ otherwise.
\end{theorem}

\begin{proof}
We apply $ev_n^{\mathcal{R}}$ to $\Omega$ coefficientwise.
Using the generating function from Lemma \ref{specialised_gen_fun_lemma},
\begin{eqnarray*}
ev_n^{\mathcal{R}}(\Omega(x_{c_1}, \ldots, x_{c_l}))ev_n^{\mathcal{R}}( \Omega(y_{c_1}, \ldots, y_{c_l})) &=& 
(1 + \sum_c cx_c)^{\otimes n} (1 + \sum_c c y_c)^{\otimes n} \\
&=& ((1 + \sum_c cx_c)(1 + \sum_c c y_c))^{\otimes n} \\
&=& (1 + \sum_c c z_c)^{\otimes n} \\
&=& ev_n^{\mathcal{R}}(\Omega(z_{c_1}, \ldots, z_{c_l})) \\
&=& ev_n^{\mathcal{R}}(\Omega(x_{c_1}, \ldots, x_{c_l})\Omega(y_{c_1}, \ldots, y_{c_l}))
\end{eqnarray*}
where $z_{c_i} = x_{c_i} + y_{c_i} + \sum_{j,k}A_{j,k}^{i} x_{c_j} y_{c_k}$ as in Proposition \ref{r_gamma_is_an_algebra_prop}. Considering the coefficients of monomials $x^{\mathbf{N}}y^{\mathbf{M}}$ shows $ev_n^{\mathcal{R}}(B_{\mathbf{N}})ev_n^{\mathcal{R}}(B_{\mathbf{M}}) = ev_n^{\mathcal{R}}(B_{\mathbf{N}}B_{\mathbf{M}})$, i.e. that $ev_n^{\mathcal{R}}$ is a homomorphism.
\end{proof}


The elements $B_{\mathbf{N}}$ can be thought of as interpolating the $b_{\mathbf{N}}$. To make this precise, we show that $ev_n^{\mathcal{R}}$ factors through $\FHG$.

\begin{corollary} \label{coeff_compatibility_cor}
We obtain a canonical homomorphism $\Psi: \RGamma \to \FHG$ via
\[
\Psi( B_{\mathbf{N}}) = \binom{t - |\mathbf{N}| + \mathbf{N}(1)}{\mathbf{N}(1)} K_{\bs \mu},
\]
where $t$ is the polynomial variable of $\mathcal{R}$, $\bs \mu(c) = (1^{\mathbf{N}(c)})$ for $c \neq 1$, and $\bs \mu(1)$ is the empty partition. This homomorphism obeys $\Phi_n \circ \Psi = ev_n^{\mathcal{R}}$.
\end{corollary}
\begin{proof}
We first show that $\Phi_n \circ \Psi = ev_n^{\mathcal{R}}$. This amounts to checking that
\[
b_{\mathbf{N}} = \binom{n - |\mathbf{N}| + \mathbf{N}(1)}{\mathbf{N}(1)} X_{\bs \mu}.
\]
Each pure tensor in $b_{\mathbf{N}}$ is an element of generalised cycle type $\bs \mu$. However, such a pure tensor has $\mathbf{N}(1)$ assigned factors equal to $1$ and $n - |\mathbf{N}|$ unassigned factors equal to $1$. So each such pure tensor arises in $\binom{n - |\mathbf{N}| + \mathbf{N}(1)}{\mathbf{N}(1)}$ ways according to how the $\mathbf{N}(1)$ assigned factors are chosen. This accounts for the multiplicative factor. To check that $\Psi$ respects multiplication, let $x,y \in \RGamma$. Because $ev_n^{\mathcal{R}}$ and $\Phi_n$ are homomorphisms,
\begin{eqnarray*}
\Phi_n(\Psi(x)\Psi(y)) &=& \Phi_n(\Psi(x))\Phi_n(\Psi(y)) \\
&=& ev_n^{\mathcal{R}}(x) ev_n^{\mathcal{R}}(y) \\
&=& ev_n^{\mathcal{R}}(xy) \\
&=& \Phi_n(\Psi(xy)).
\end{eqnarray*}
We now use Proposition \ref{specialisation_distinguishing_prop} to conclude that $\Psi(x)\Psi(y) = \Psi(xy)$.
\end{proof}
